\subsection{循环模的直和}

设 A 为一个交换环。回顾（II，第 220 页，命题 22），一个循环 A-模同构于商模 A/a，其中 a 是 A 的一个理想。我们将在后面（第 4 节）看到，主理想整环上的每个有限生成模都是有限多个循环模的直和。

命题 1. --- 设 E 为交换环 A 上的一个模；假设 E 是 n 个循环模 \(A/\mathfrak{a}_{k}\) ( \(1 \leqslant k \leqslant n\) ) 的直和，其中 \(a_{k}\) 是 A 的理想；

那么，对于每个整数 \(p > 0\) ，A-模 \(\bigwedge^p \mathbf{E}\) 同构于诸模 \(\mathbf{A} / \mathfrak{a}_{\mathrm{H}}\) 的直和，其中对于每个 \(p\) -元素子集 \(\mathbf{H} = \{k_1, \dots, k_p\}\) （ \([1, n]\) 的），理想 \(\mathfrak{a}_{\mathrm{H}}\) 是 \(\sum_{j=1}^{p} \mathfrak{a}_{k_j}\) .

设 \(x_{k}\) 为 \(A / a_{k}\) 的典范生成元，即 A 的单位元的像，因此 E 是诸 \(Ax_{i}\) ( \(1 \leqslant i \leqslant n\) ) 的直和。于是我们知道（III，第 515 页，命题 10），外代数 \(\bigwedge_{n} E\) 作为 A-模同构于张量积 \(\otimes_{i=1}^{n} (\bigwedge(Ax_{i}))\) 。现在 \(\bigwedge(Ax_{i})\) 正是直和 \(A \oplus Ax_{i}\) ，因为

\[
\bigwedge^ {p}
\]

\(Ax_{i}\) 中任意两个元素的外积为零，因此 \(\Lambda E\) 是诸模 \(\mathbf{M}_{\mathrm{H}} = (\mathbf{A}x_{k_{1}}) \otimes \ldots \otimes (\mathbf{A}x_{k_{p}})\) 的直和，其中 \(H = \{k_{1}, \ldots, k_{p}\}\) 取遍 {[}1, n{]} 的 p 元子集（ \(k_{1} < \ldots < k_{p}\) ）；现在已知 \(M_{H}\) 同构于 \(A/a_{H}\) （II，第 257 页，推论 4），这就完成了证明。

我们现在看到，在命题 1 的记号下，如果诸理想 \(\mathfrak{a}_k\) 构成一个递增序列，则它们完全由模 E 确定。更精确地说：

命题 2. --- 设 A 是一个交换环，且 E 是 n 个循环模 \(A/a_{k}\) 的直和，其中 \(a_{k}\) 满足 \(a_{1} \subset a_{2} \subset \ldots \subset a_{n}\) 。那么，对于 \(1 \leqslant p \leqslant n\) ，理想 \(a_{p}\) 是 \(\bigwedge^{p} E\) 的零化子；如果 \(a_{n} \neq A\) ，则 \(\bigwedge^{p} E \neq 0\) （对 \(1 \leqslant p \leqslant n\) ），且 \(\bigwedge^{m} E = 0\) （对 m \textgreater{} n）.

确实，在命题 1 的记号中，我们有 \(\mathfrak{a}_{\mathrm{H}} = \mathfrak{a}_{s(\mathrm{H})}\) ，其中 \(s(\mathrm{H})\) 表示子集 H 的最大元素。由于 \(s(\mathrm{H}) \geqslant p\) 对每个 \(p\) -元素子集 H 成立，且 \(s(\mathrm{H}) = p\) 对 \(\mathrm{H} = \{1, \dots, p\}\) 成立，由此可得， \(\mathfrak{a}_p\) 是诸 \(\mathfrak{a}_{\mathrm{H}}\) 的交，其中 H 取遍 \(p\) -元素子集（ \([1, n]\) 的）；因此理想 \(\mathfrak{a}_p\) 确实是 \(\bigwedge^p E\) 的零化子（由命题 1）。

推论. --- 在命题 2 的记号下，若 \(\mathfrak{a}_n \neq \mathbf{A}\) ，且 E 也同构于 \(m\) 个循环模 \(\mathbf{A} / \mathfrak{a}_j'\) 的直和，满足 \(\mathfrak{a}_1' \subset \mathfrak{a}_2' \subset \ldots \subset \mathfrak{a}_m' \neq \mathbf{A}\) ，则 \(m = n\) ，且 \(\mathfrak{a}_k = \mathfrak{a}_k'\) 对 \(1 \leqslant k \leqslant n\) 成立（«诸 \(\mathfrak{a}_k\) 的唯一性»).
% semantic-fragments: legacy-input-seam
\relax{}
